\input format.tex

%\usepackage{arydshln}
\usepackage{graphicx}
\graphicspath{{cores/}}

\begin{document}
\color{gray2}

\vspace*{0mm}
\fontsize{9.3pt}{9.8pt}\selectfont
\begin{center}
{\bf\sanhao 健康建议}
\end{center}

\begin{tctabularx}{tabularx={m{15.4cm}}}
\\[-6.769pt]
\cellcolor{kelly} \raisebox{1pt}{\color{black70}\xiaosihao\lantxh\bf 肠道健康总说明} \\
\end{tctabularx}

{
\lantxh
\vspace*{-4.25mm}
\fontsize{8.8pt}{11pt}\selectfont
\begin{longtable}{|@{~~~}m{4cm}<{\centering}|m{10.96cm}|}
\hline
\parbox[c]{\hsize}{\vskip6pt\raisebox{-4ex}{\includegraphics[width=1.8cm]{INSUFFICIENCY-yellow.pdf}}~~\xiaowuhao \color{gray2} 消化和吸收\vskip6pt} &
\parbox[c]{\hsize}{\vskip6pt\xiaowuhao\color{gray2} {
通过菌群测得您的
蛋白质消化吸收能力正常、脂肪消化吸收能力低、碳水化合物消化吸收能力正常；15种主要代谢物中
牛磺酸代谢、鞘脂代谢、亚油酸代谢等指标优于参考人群；脂肪总量、甘油三酯代谢、甘油磷脂代谢指标存在异常。
}\vskip4pt}\\
\hline
\parbox[c]{\hsize}{\vskip6pt\raisebox{-4ex}{\includegraphics[width=1.8cm]{INFLAMMATION-yellow.pdf}}~~\xiaowuhao\color{gray2} 炎症和免疫\vskip6pt} &
\parbox[c]{\hsize}{\vskip6pt\xiaowuhao\color{gray2} {
您的
肠道净化功能正常；肠道排毒功能正常；肠道免疫功能较差。}\vskip4pt}\\
\hline
\parbox[c]{\hsize}{\vskip6pt\raisebox{-4ex}{\includegraphics[width=1.8cm]{MICROBIOTA-low.pdf}}~~\xiaowuhao\color{gray2} 肠道菌群\vskip6pt} &
\parbox[c]{\hsize}{\vskip6pt\xiaowuhao\color{gray2} {
您的肠道菌群以拟杆菌属为主。肠道菌群多样性水平
较差，菌群失调的风险较高。
有2类细菌指标与参考人群相比存在异常。
}\vskip4pt}\\
\hline
\parbox[c]{\hsize}{\vskip6pt\raisebox{-4ex}{\includegraphics[width=1.8cm]{INFECTION-yellow.pdf}}~~\xiaowuhao\color{gray2} 致病菌\vskip6pt} &
\parbox[c]{\hsize}{\vskip6pt\xiaowuhao\color{gray2} {
%您的肠道中检测到1种可能导致疾病的致病菌，其中种食源性致病菌，种机会致病菌，含量未超标，但在您免疫能力低下时仍可能会损害您的肠道健康。
您的肠道中检测到1种可能导致疾病的致病菌，含量未超标，但在您免疫能力低下时仍可能会损害您的肠道健康。
}\vskip4pt}\\
\hline
\parbox[c]{\hsize}{\vskip6pt\raisebox{-4ex}{\includegraphics[width=1.8cm]{IMBALANCE-yellow.pdf}}~~\xiaowuhao\color{gray2} 代谢平衡\vskip6pt} &
\parbox[c]{\hsize}{\vskip6pt\xiaowuhao\color{gray2} {
您的代谢基本平衡，11种重要有机小分子代谢中胆碱代谢异常。
}\vskip4pt}\\
\hline
\end{longtable}
}

\fontsize{9.3pt}{9.8pt}\selectfont
\begin{tctabularx}{tabularx={m{15.4cm}}}
\\[-6.769pt]
\cellcolor{kelly} \raisebox{1pt}{\color{black70}\xiaosihao\lantxh\bf GIhealth 提示您} \\
\end{tctabularx}

{
\lantxh
\vspace*{-4.25mm}
\fontsize{8.8pt}{11pt}\selectfont
\begin{longtable}{|m{15.45cm}|}
\hline
\parbox[c]{\hsize}{\vskip11pt\begin{compactenum}[\mbox{1、}]
{\xiaowuhao\color{gray2}\item[1、]
您的肠道菌群多样性
较差，仅高于6.10{\%}的参考人群。肠道有益菌含量
较低，仅高于6.43{\%}的参考人群，有害菌含量
正常，高于26.63{\%}的参考人群。
检出1种致病菌，含量在正常范围内。
}
{\xiaowuhao\color{gray2}\item[2、]
%通过对您的肠道菌群分析，您患2型糖尿病、结直肠癌、炎症性肠病的风险均偏高。
通过对您的肠道菌群分析，您的肠道菌群可能增加患2型糖尿病、结直肠癌、炎症性肠病的风险。
}
{\xiaowuhao\color{gray2}\item[3、]
您肠道内参与胆碱代谢的菌群指标与参考人群相比存在异常。代谢基本平衡，请关注您的异常指标，定期监测肠道健康。
}
{\xiaowuhao\color{gray2}\item[4、]
您肠道内参与脂肪总量、甘油三酯代谢、甘油磷脂代谢的菌群指标与参考人群相比存在异常，请关注您的异常指标，定期监测肠道健康。
}
{\xiaowuhao\color{gray2}\item[5、]}
\end{compactenum}\vskip2pt}\\
\hline
\end{longtable}
}

\end{document}
